%% sections/01_introduction.tex
\section{Introduction}
\label{sec:introduction}

Distributed machine learning on relational networks, structured as Federated Graph Neural Networks (FedGNNs), has emerged as a cornerstone paradigm for privacy-preserving collaborative intelligence across regulated domains, including anti-money laundering in financial consortia, legal recidivism risk assessment, and clinical healthcare networks~\citep{mcmahan2017communication, kipf2016semi, velickovic2017graph, kairouz2021advances}. In these cross-silo enterprise federations, graph structures encode sensitive relationships between legal entities, bank accounts, or patient cohorts. Legal mandates, such as the European Union AI Act, the General Data Protection Regulation (GDPR), and the U.S. Equal Credit Opportunity Act, strictly forbid centralizing raw relational data across sovereign institutional boundaries. Concurrently, these statutes demand that predictive algorithms ensure non-discriminatory group demographic fairness while maintaining reliable operation in the presence of malicious or compromised participant nodes.

\paragraph{The Trustworthy Federated Learning Trilemma}
A fundamental architectural challenge in decentralized graph learning is that the system must reconcile three inherently competing objectives:
\begin{enumerate}[leftmargin=*]
    \item \textbf{Unconstrained Task Utility ($U$):} Maximizing predictive accuracy by solving $\min_\theta \mathcal{L}_{\mathrm{task}}(\theta)$. Standard single-objective aggregation protocols (e.g., FedAvg, FLTrust) dedicate the entirety of model parameter updates to minimizing empirical task loss, reaching high accuracy on clean data. However, because they allocate zero gradient capacity to algorithmic fairness, they leave demographic disparity unconstrained ($\dpd \ge 0.055$ on German Credit) and remain blind to adversarial perturbations injected into the fairness subspace.
    \item \textbf{Group Demographic Fairness ($F$):} Enforcing demographic parity ($\dpd(\theta) \le \tau$) across protected populations. Existing fairness-aware federated protocols (e.g., FairFed, $F^2\text{GNN}$, FedGraph-Fair, BFWA) rely on client-reported scalar disparity metrics to coordinate aggregation weights.
    \item \textbf{Byzantine Robustness ($R$):} Bounding the influence of arbitrary or malicious clients ($w_{\mathrm{adv}} \le 1/K$). Standard robust geometric filters (e.g., Coordinate Median, Krum, FLAME) assume symmetric parameter distributions, treating outlier vectors as anomalies.
\end{enumerate}
In decentralized relational learning, these three objectives form the \emph{Trustworthy FL Trilemma}: unconstrained optimization along any one or two axes inevitably induces structural failure on the remaining dimensions.

\paragraph{Failure Mode 1: The Self-Reporting Metadata Trap}
The dominant design pattern in contemporary fairness-aware federated learning coordinates client aggregation weights based on self-reported scalar telemetry, such as local demographic disparity $\widehat{\dpd}_k$ or local empirical loss~\citep{ezzeldin2023fairfed, zhou2026fairgfl, f2gnn2023, khan2026fedgraphfair}, including our own earlier conference formulation~\citep{dang2026fedfairgnn}. In adversarial environments governed by Kerckhoffs' principle, this reliance creates a critical attack surface. As established in our empirical vulnerability evaluation (Section~\ref{subsec:vulnerability}, Table~\ref{tab:metadata_capture}), Byzantine adversaries can transmit poisoned gradient updates while falsifying local disparity claims ($\widehat{\dpd}_k \approx 0.000$). Across six distinct self-reporting aggregation rules, adversaries systematically hijack the weighting mechanism: in $F^2\text{GNN}$, the attacker captures $w_{\mathrm{adv}} = 0.9737$ of global model authority ($4.87\times$ its uniform share), inflating downstream disparity by $\Delta\dpd = +0.4638$ ($p = 1.9 \times 10^{-9}$). Client-reported metadata interfaces are fundamentally fragile under adversarial incentives.

\paragraph{Failure Mode 2: The Information-Theoretic Barrier of Local Differential Privacy}
A natural defense is to mask self-reported scalar metrics via Local Differential Privacy (LDP). However, applying noise to disparity statistics encounters two insurmountable theoretical barriers:
\begin{enumerate}[leftmargin=*]
    \item \emph{Perpetual Expectation Bias (Lemma~\ref{lem:folded_normal}):} Because disparity is evaluated via an absolute difference operator ($\widetilde{\dpd}_k = |\tilde\mu_{k,0} - \tilde\mu_{k,1}| \ge 0$), zero-mean Gaussian noise is transformed into a central folded-normal distribution~\citep{leone1961folded}. Conditioned on past rounds $\mathcal{F}_{t-1}$, the expectation remains strictly positive ($\EE[\widetilde{\dpd}_k] = \tilde\sigma_k \sqrt{2/\pi} > 0$), and almost surely $\sum_k w_k \widetilde{\dpd}_k \ge \min_k \widetilde{\dpd}_k > 0$. The noise cannot cancel across silos, dilating the apparent disparity budget by up to $+10{,}891\%$ (Table~\ref{tab:ldp_barrier}).
    \item \emph{Minimax Observability Barrier (Theorem~\ref{thm:minimax_lecam}):} Applying Le Cam's two-point method and Pinsker's inequality, we prove that distinguishing fair from unfair states under LDP is bounded from below by minimax testing risk $\Rstar \ge 43.5\%$ at privacy budget $\epsilon = 8.0$ (rising to $\ge 49.0\%$ at $\epsilon = 1.0$, Table~\ref{tab:ldp_barrier}). By the Data Processing Inequality, no server-side post-processing can lower this observational floor.
\end{enumerate}
Furthermore, while applying DP-SGD across the full high-dimensional parameter tensor reduces membership leakage, it destroys predictive utility, incurring a parameter degradation penalty $115\times$ more severe than scalar-channel privacy (Figure~\ref{fig:privacy_bail}).

\paragraph{Failure Mode 3: Geometric Defense Failures and Enterprise Constraints}
Turning to classical Byzantine defenses reveals equally severe structural vulnerabilities:
\begin{enumerate}[leftmargin=*]
    \item \emph{Fairness Subspace Blindness (Theorem~\ref{thm:subspace_blindness}):} Anchor-free geometric detectors (e.g., Krum, FLAME) operate over the cosine geometry of task parameter updates. Because task classification coordinates dominate the gradient spectrum by orders of magnitude relative to low-rank fairness parameters ($\mathcal{D}_{\mathrm{cos}} \approx \mathcal{D}_{\mathrm{cos}}^{\mathrm{task}}$), proximity-stealth adversaries can inject substantial fairness degradation while remaining geometrically indistinguishable inside the benign clustering ball ($\rho = 0.85$).
    \item \emph{The Shielding Paradox (Median Backfire):} On non-IID graph federations with modular community partitions, benign clients in minority topological clusters naturally exhibit elevated gradient variance. Symmetric coordinate-wise median filtering systematically purges benign minority gradients before eliminating stealth adversaries. As demonstrated in Section~\ref{subsec:stress_testing} (Table~\ref{tab:adaptive_poisoner}), median filtering backfires, inflating adversary aggregation share by $+82.9\%$ relative to canonical averaging ($w_{\mathrm{adv}} = 0.5343$ at $f/K=0.4$).
\end{enumerate}
Finally, heavyweight cryptographic mechanisms (e.g., ZK-SNARKs, Secure MPC) and hardware enclaves (TEEs) are practically prohibitive in cross-silo graph federations due to extreme graph message-passing computational overhead and heterogeneous enterprise infrastructure.

\paragraph{The \MethodName{} Paradigm: Exogenous Server Anchoring}
Eliminating these flawed avenues forces a necessary conclusion: robust, fair aggregation cannot depend on client-reported metadata or symmetric unanchored geometry. Instead, the server must evaluate incoming parameter deltas against an independent, physically isolated directional reference. 

We propose \MethodName{}, an auditable cross-silo framework that achieves balanced Pareto operation across the Trustworthy FL Trilemma. \MethodName{} implements a principled two-tier decoupling:
\begin{itemize}[leftmargin=*]
    \item \textbf{Client Side (FTGD):} Local silos execute Fair Targeted Gradient Descent with first-order orthogonal projection (Theorem~\ref{thm:ftgd_orthogonality}), filtering sensitive demographic gradients locally prior to parameter transmission.
    \item \textbf{Server Side (Canonical FU-Alignment):} The server maintains an isolated audit anchor dataset $\Droot$ (consisting of $125$ stratified nodes) to compute an exogenous dual-objective target gradient $\gtarg = \gtask^{\mathrm{srv}} + \alpha \gfair^{\mathrm{srv}}$. Canonical aggregation combines rectified directional cosine gating, dynamic norm rescaling, and exponential moving average (EMA) smoothing. This mathematical formulation achieves exact syntactic metadata immunity ($\Delta w = 0.0000$, Theorem~\ref{thm:syntactic_immunity}) while dampening round-to-round weight turbulence. To mitigate crude unbounded scaling attacks ($c = 100$) without inducing the shielding paradox under stealth, a defense-in-depth coordinate median screening variant (\texttt{robust\_fu\_shapley}) is provided as a static configuration arm.
\end{itemize}

\paragraph{Key Scientific Contributions}
This work delivers four primary contributions to trustworthy decentralized graph learning:
\begin{enumerate}[leftmargin=*]
    \item \textbf{Formalization of the Trustworthy FL Trilemma and Threat Hierarchy:} We formalize the structural conflict between Task Utility ($U$), Demographic Fairness ($F$), and Byzantine Robustness ($R$) in decentralized graph learning. Under Kerckhoffs' principle, we formulate a 3-tier adversarial hierarchy and prove exact syntactic metadata immunity ($\Delta w = 0.0000$) as an architectural invariant (Theorem~\ref{thm:syntactic_immunity}), categorizing aggregation rules into naive scalar consumers, endogenous references, and exogenous anchors.
    \item \textbf{Information-Theoretic Limits and Subspace Blindness:} We prove the perpetual expectation bias of folded-normal perturbations under LDP (Lemma~\ref{lem:folded_normal}) and establish the Le Cam minimax observability lower bound $\Rstar \ge 43.5\%$ at $\epsilon = 8.0$ (rising to $\ge 49.0\%$ at $\epsilon = 1.0$, Theorem~\ref{thm:minimax_lecam}, Table~\ref{tab:ldp_barrier}). Furthermore, we prove Fairness Subspace Blindness for anchor-free geometric detectors (Theorem~\ref{thm:subspace_blindness}), explaining why clustering algorithms fail against proximity-stealth fairness attacks.
    \item \textbf{Two-Tier Architecture and Formal Axiomatic Verification:} We propose \MethodName{}, integrating local FTGD orthogonal projection with server-side Canonical FU-Alignment. To ensure strict formal correctness, Theorem~\ref{thm:ftgd_orthogonality} (first-order orthogonality) and Propositions~\ref{prop:linear_shapley}--\ref{prop:null_player} (first-order linear decomposition in $O(KP)$, simplex validity, and null-player guarantee) are formally verified with machine-checked proof certificates.
    \item \textbf{Empirical Verification and Auditable Governance:} Across four real-world graph benchmarks (Pokec-z, German Credit, Bail Recidivism, Credit Default), 14 baselines, and three attack tiers, we demonstrate that \MethodName{} establishes an intentional, balanced operating point on the Pareto frontier (trading $2.00\%$ AUC margin for a $41.5\%$ disparity reduction on Pokec-z, while achieving statistical equivalence on benign data: $\Delta\mathrm{AUC} = +0.0015, p = 0.625$). We provide an $O(KP)$ per-round governance ledger ($r = 0.7662 < 0.80$ against exponential combinatorial Shapley), and transparently report operational boundaries under extreme Dirichlet non-IID shifts and white-box memory compromises.
\end{enumerate}

\paragraph{Paper Organization}
The remainder of this manuscript is organized as follows. Section~\ref{sec:related_work} positions our framework within five foundational literature paradigms and outlines operational boundaries. Section~\ref{sec:methodology} presents the formal problem formulation, the \MethodName{} architecture, theoretical observability bounds, and formal derivations of governance properties. Section~\ref{sec:experiments} reports extensive empirical stress tests, SOTA benchmark comparisons, ablation studies, and causal attribution analyses. Section~\ref{sec:conclusion} concludes the paper with operational deployment guidelines and research limitations.
